\section{算法设计}

NPQR 是一个神经辅助的组合式量子路由算法。它的核心不是单独的神经网络模型，而是
“神经网络评分 + 经典搜索控制”的整体流程。图 \ref{fig:npqr-flow} 展示了主要
步骤，算法 \ref{alg:npqr} 给出高层伪代码。

\begin{figure}[H]
  \centering
  \begin{tikzpicture}[
  font=\small,
  phasebox/.style={
    rectangle,
    rounded corners=2pt,
    draw=blue!55!black,
    fill=blue!6,
    align=center,
    minimum width=2.6cm,
    minimum height=0.78cm
  },
  verify/.style={
    rectangle,
    rounded corners=2pt,
    draw=green!45!black,
    fill=green!8,
    align=center,
    minimum width=2.6cm,
    minimum height=0.78cm
  },
  arrow/.style={-{Latex[length=2mm]}, thick, draw=gray!70!black}
]
  \node[phasebox] (input)   at (0,0)       {线路与拓扑};
  \node[phasebox] (pre)     at (3.4,0)     {图预处理};
  \node[phasebox] (mapping) at (6.8,0)     {初始映射};

  \node[phasebox] (beam)    at (6.8,-1.65) {束搜索};
  \node[phasebox] (score)   at (3.4,-1.65) {动作评分};
  \node[phasebox] (front)   at (0,-1.65)   {前沿推进};

  \node[phasebox] (repair)  at (0,-3.30)   {局部修复};
  \node[verify]   (replay)  at (3.4,-3.30) {复放验证};
  \node[verify]   (output)  at (6.8,-3.30) {线路与指标};

  \draw[arrow] (input) -- (pre);
  \draw[arrow] (pre) -- (mapping);
  \draw[arrow] (mapping) -- (beam);
  \draw[arrow] (beam) -- (score);
  \draw[arrow] (score) -- (front);
  \draw[arrow] (front) -- (repair);
  \draw[arrow] (repair) -- (replay);
  \draw[arrow] (replay) -- (output);
\end{tikzpicture}

  \caption{NPQR 神经辅助量子路由流程}
  \label{fig:npqr-flow}
\end{figure}

从职责划分看，NPQR 的每个模块都解决搜索中的一个具体问题。表
\ref{tab:component-roles} 概括了这些模块的输入、输出和作用。

\begin{table}[H]
  \centering
  \caption{组件职责}
  \label{tab:component-roles}
  \begin{tabularx}{\linewidth}{lYYY}
    \toprule
    模块 & 输入 & 输出 & 作用 \\
    \midrule
    依赖关系 & 门序列 & 前沿门、前驱关系 & 判断哪些门可以在当前状态推进。\\
    拓扑预处理 & 硬件耦合图 & 距离矩阵、边集合 & 为动作评分和合法性检查提供图特征。\\
    初始映射 & 逻辑交互图、物理拓扑 & 候选初始布局 & 降低后续 SWAP 搜索压力。\\
    动作筛选 & 当前映射、前沿门 & 候选 SWAP 集合 & 删除弱相关动作，控制分支数量。\\
    模型评分 & 状态特征、候选动作 & 动作偏好分数 & 为搜索提供学习得到的局部判断。\\
    束搜索 & 候选状态、累计得分 & 多条候选路线 & 保留多条高分路线，提升搜索质量。\\
    局部修复 & 接近完成的路线 & 修复后的尾部轨迹 & 处理少量剩余门停滞的情况。\\
    轨迹复放 & 执行轨迹、硬件拓扑 & 合法性检查结果 & 确认最终线路满足拓扑约束。\\
    \bottomrule
  \end{tabularx}
\end{table}


\subsection{图预处理与前沿门}

系统首先解析输入电路，提取门序列、量子比特集合和双比特门依赖关系。随后将目标
芯片表示为无向图，并计算物理量子比特之间的最短路径距离矩阵。该预处理把后续
多次距离查询变成常数时间访问，是典型的时空权衡。

在路由过程中，系统维护当前前沿门集合。若前沿双比特门在当前映射下已经相邻，
系统直接执行该门并推进门依赖；若不能执行，系统生成候选 SWAP 动作，并进入评分
与搜索流程。

\subsection{初始映射选择}

初始映射决定逻辑量子比特最开始放置在哪些物理量子比特上。若初始映射与电路的
交互结构匹配，后续需要的 SWAP 会减少。NPQR 会根据逻辑交互关系和物理拓扑生成
多个候选映射，并对候选进行排序。该步骤把原始电路路由问题转化为较好的搜索起点，
降低后续局部搜索的压力。

\subsection{神经网络动作评分}

在每个非终止状态下，神经网络根据当前映射、前沿门、物理距离和候选 SWAP 动作
输出动作偏好分数。该分数不直接决定完整答案，而是作为束搜索中的一个启发式信号。
这种设计使模型能够参与决策，同时避免模型一次性生成完整电路带来的不可控风险。

为了便于解释，可以把一个候选动作 $a$ 在状态 $s$ 下的综合评分写成
\[
  F(s,a)=\lambda_1 P_\theta(a\mid s)-\lambda_2 H_{\mathrm{frontier}}(s,a)
  -\lambda_3 C_{\mathrm{depth}}(s,a)-\lambda_4 C_{\mathrm{repeat}}(s,a),
\]
其中 $P_\theta(a\mid s)$ 表示神经网络给出的动作偏好，$H_{\mathrm{frontier}}$
表示前沿门距离启发式代价，$C_{\mathrm{depth}}$ 表示对深度增长的惩罚，
$C_{\mathrm{repeat}}$ 表示对重复无效移动的惩罚。实际系统可以使用等价的排序
逻辑；该公式强调的是：神经网络提供主要偏好，但最终决策仍受拓扑距离和搜索状态
共同约束。

\subsection{有界束搜索}

单步贪心容易因为局部最优而卡住。NPQR 使用有界束搜索保留多条候选路线。每次
展开时，系统对候选状态进行评分，只保留排名靠前的若干条路线继续搜索。束宽越大，
搜索质量通常越高，但时间和空间消耗也随之增加。

束搜索的关键在于把“当前最优动作”扩展为“当前较优路线集合”。设第 $t$ 层保留的
状态集合为 $\mathcal{B}_t$，系统会对每个 $s\in \mathcal{B}_t$ 展开若干动作并生成
候选集合 $\mathcal{C}_{t+1}$，然后按累计评分取前 $B$ 个状态：
\[
  \mathcal{B}_{t+1}=\mathrm{TopB}(\mathcal{C}_{t+1}).
\]
这样做牺牲了一部分运行时间，但能减少单步评分误差造成的不可逆失败。

\subsection{触发剪枝与局部修复}

对普通样例，过重搜索会增加不必要耗时；对交互密集或随机性强的样例，过轻搜索又
可能漏掉可行路线。NPQR 使用前沿触发规则识别更难的局部结构，只在必要时启用更强
的前沿搜索。同时，动作剪枝会删除明显较差的 SWAP 候选，控制搜索树分支数。

触发规则主要关注前沿门的拥塞程度、候选动作是否能接触关键前沿量子比特、以及
普通展开是否出现停滞。若状态表现为“前沿门距离长期不能下降”或“候选动作虽然得分
相近但影响范围不同”，系统会扩大对前沿相关动作的考虑范围。动作剪枝则保留与当前
前沿门、短路径和神经评分更相关的候选，避免在所有硬件边上无差别展开。

当一条路线已经完成大部分门但在末端卡住时，NPQR 会启用有界后缀修复。后缀修复
不是全局穷举，而是在有限深度内围绕剩余门进行局部搜索，用于补齐接近完成的路线。

\subsection{轨迹复放验证}

最终结果必须通过轨迹复放。系统会从初始映射开始重新模拟所有门和 SWAP，检查每
一步是否满足硬件边约束，并确认所有原始门都已执行。只有通过复放验证的路线才会
作为 NPQR 结果返回。该步骤把搜索输出转化为可审计的工程结果。

轨迹复放还用于区分“搜索过程看似完成”和“输出结果确实可执行”。如果某一步双比特
门不在硬件边上，或某个原始门没有被执行，复放都会失败。这个设计对课程项目很重要：
它把算法输出和硬件约束重新绑定，避免只报告一个未经验证的启发式分数。

\begin{figure}[H]
  \centering
  \begin{minipage}{0.94\linewidth}
  \begin{tabbing}
  \hspace{1.8em}\=\hspace{1.8em}\=\hspace{1.8em}\=\kill
  \textbf{算法 1：NPQR 神经辅助量子路由}\\
  \textbf{输入：}量子电路 $C$，硬件图 $G$，束宽 $B$，修复深度 $R$\\
  \textbf{输出：}路由后电路、SWAP 数、轨迹或失败标记\\
  1. 解析 $C$，构建门依赖和前沿门集合；\\
  2. 计算 $G$ 上的距离矩阵，生成初始映射候选；\\
  3. 将候选映射转化为初始搜索状态集合；\\
  4. \textbf{while} 存在未完成状态 \textbf{do}\\
  \> 4.1 执行所有当前可直接执行的前沿门；\\
  \> 4.2 若状态完成，进行轨迹复放并返回结果；\\
  \> 4.3 生成合法 SWAP 候选动作；\\
  \> 4.4 用神经网络和距离启发式为动作评分；\\
  \> 4.5 根据触发规则决定是否扩大前沿搜索；\\
  \> 4.6 剪枝弱动作，展开候选状态；\\
  \> 4.7 保留评分最高的 $B$ 个状态；\\
  \> 4.8 若状态接近完成但卡住，尝试深度不超过 $R$ 的后缀修复；\\
  5. \textbf{return} 失败标记。
  \end{tabbing}
  \end{minipage}
  \caption{NPQR 高层伪代码}
  \label{alg:npqr}
\end{figure}
